% Faithful formal-publication transcription: complete Section5.
\section*{5 Proof of Multileaf Lemma}
\begin{authorclaim}{Multileaf Lemma (restated)}
For every partial $g$ and $0<\epsilon\lll\delta\ll1$, there exists a distribution $\mu=\frac12(\mu^0+\mu^1)$ over $\operatorname{dom}(g)$ and a hard side $b\in\{0,1\}$ such that for every deterministic tree $t$ taking inputs from $\operatorname{dom}(g^n)$ and having $\operatorname{depth}(t)/(n\overline R_\epsilon(g))\lll\delta$,
\[\forall y\in\{0,1\}^n:\quad\Pr_{x\sim\mu^y}[\ell_x^t\text{ is }(\delta,\epsilon,\mu,b)\text{-noisy}]\ge1-\delta.\]
\end{authorclaim}
\begin{proof}
Apply Leaf Lemma with parameters $\epsilon$ and $\dot\delta:=\delta^3$ (assuming suitably $0<\epsilon\lll\dot\delta\ll1$) to obtain $\mu=\frac12(\mu^0+\mu^1)$ and $b\in\{0,1\}$ that satisfy the lemma for trees of depth at most $\dot\delta^c\overline R_\epsilon(g)$. Fix $y\in\{0,1\}^n$ and a deterministic tree $t$ over $\operatorname{dom}(g^n)$ with $\operatorname{depth}(t)\le\dot\delta^{c+4}n\overline R_\epsilon(g)$ (where we chose $3(c+4)$ as the exponent hidden by $\lll$).
Here is the plan for our proof. An input $x\in\operatorname{dom}(g^n)$ can be seen as inducing several subtrees of $t$ corresponding to distinct coordinates $i\in[n]$. Indeed, define $t_{x,i}$ as the tree over inputs from $\operatorname{dom}(g)$ that is obtained from $t$ by substituting $x$ as its input variables except retaining $x_i$ as free variables. If we can show that $t_{x,i}$ has shallow depth in expectation over an input $z\sim\mu^{y_i}$ then we can hope to use and argue that the reached leaf $\ell_z\in L(t_{x,i})$ (which is one of the $n$ components of a leaf of $t$) is typically $(\epsilon,\mu,b)$-noisy.
Let us formalise this plan. Let $x\sim\mu^y$ henceforth. For $i\in[n]$ we define two events
\begin{align*}
\text{$i$-th tree is shallow: }\quad S_i(x)&\Longleftrightarrow\mathbb E_{z\sim\mu^{y_i}}[q(t_{x,i},z)]\le\dot\delta^c\overline R_\epsilon(g),\\
\text{$i$-th leaf is noisy: }\quad N_i(x)&\Longleftrightarrow\ell_{x_i}\in L(t_{x,i})\text{ is }(\epsilon,\mu,b)\text{-noisy}.
\end{align*}
Note that Leaf Lemma states $\Pr_x[N_i\mid S_i]\ge1-\dot\delta$. Thinking of $S_i$ and $N_i$ as indicator variables, we define $S:=\frac1n\sum_iS_i$ and $N:=\frac1n\sum_iN_i$. With this notation, Multileaf Lemma becomes equivalent to
\begin{equation}\tag{9}\Pr_x[N\ge1-\delta]\ge1-\delta.\end{equation}
To show this, we compute as follows (using Claim13 that is proved below)
\begin{align*}
\mathbb E_x[N]&=\frac1n\sum_i\Pr_x[N_i]\ge\frac1n\sum_i(1-\dot\delta)\Pr[S_i]\quad\text{(Leaf Lemma)}\\
&=(1-\dot\delta)\mathbb E_x[S]\ge(1-\dot\delta)(1-\dot\delta)\quad\text{(Claim13)}\\
&\ge1-\delta^2.\quad(\dot\delta:=\delta^3\ll1)
\end{align*}
Hence(9) follows by applying Markov's inequality to the nonnegative random variable $1-N\ge0$. This completes the proof apart from the following claim.
\end{proof}
\begin{authorclaim}{Claim13}
$\mathbb E_x[S]\ge1-\dot\delta$.
\end{authorclaim}
\begin{proof}
Let $q_i(t,x)$ denote the number of queries made by $t$ to the $i$-th component of $x$. Define $x^{i\leftarrow z}$ as a copy of $x$ but where $z$ is inserted at the $i$-th component. Note that $q_i(t,x^{i\leftarrow z})=q(t_{x,i},z)$. Linearity of expectation gives
\begin{equation}\tag{10}\sum_{i\in[n]}\mathbb E_x[q_i(t,x)]\le\operatorname{depth}(t)\le\dot\delta^{c+4}n\overline R_\epsilon(g).\end{equation}
Define $I\subseteq[n]$ as the set of coordinates $i$ satisfying
\begin{equation}\tag{11}\mathbb E_x[q_i(t,x)]\le\dot\delta^{c+2}\overline R_\epsilon(g).\end{equation}
We have that $|I|\ge(1-\dot\delta^2)n$ as otherwise more than $\dot\delta^2n$ terms in the sum(10) are larger than $\dot\delta^{c+2}\overline R_\epsilon(g)$ contradicting the upper bound on $\operatorname{depth}(t)$. Fix $i\in I$. Sampling $x\sim\mu^y$ is equivalent to first taking $x\sim\mu^y$, then sampling independently $z\sim\mu^{y_i}$, and finally outputting $x^{i\leftarrow z}$. Hence
\[\mathbb E_x\mathbb E_{z\sim\mu^{y_i}}[q_i(t,x^{i\leftarrow z})]=\mathbb E_x[q_i(t,x)]\le\dot\delta^{c+2}\overline R_\epsilon(g).\]
We get from Markov's inequality and the above that
\begin{equation}\tag{12}\Pr_x[\neg S_i]=\Pr_x\left[\mathbb E_{z\sim\mu^{y_i}}[q_i(t,x^{i\leftarrow z})]>\dot\delta^c\overline R_\epsilon(g)\right]\le\dot\delta^2.\end{equation}
In conclusion,
\[\mathbb E_x[S]\ge\frac1n\sum_{i\in I}\Pr_x[S_i]\ge\frac1n|I|\cdot(1-\dot\delta^2)\ge(1-\dot\delta^2)^2\ge1-\dot\delta.\]
\end{proof}
